\documentclass[10pt,a4paper]{article}

%----------------------------------------------------------------------------------------
%	PACKAGES
%----------------------------------------------------------------------------------------
\usepackage[left=0.6in,top=0.6in,right=0.6in,bottom=0.6in]{geometry}
\usepackage[utf8]{inputenc}
\usepackage[T1]{fontenc}
\usepackage[default]{sourcesanspro}
\usepackage{xcolor}
\usepackage{titlesec}
\usepackage{enumitem}
\usepackage{parskip}
\usepackage{hyperref}
\usepackage{fontawesome5}

%----------------------------------------------------------------------------------------
%	COLORS & STYLES
%----------------------------------------------------------------------------------------
\definecolor{titleblue}{HTML}{00199e}
\definecolor{subtitleblue}{HTML}{2ec1e0}
\definecolor{darktext}{HTML}{222222}
\definecolor{linkblue}{HTML}{0B5FA5}
\definecolor{subtle}{HTML}{666666}

\color{darktext}
\linespread{0.9}
\setlength{\parskip}{0.08em}

\titleformat{\section}{\normalsize\bfseries\color{titleblue}}{}{0em}{}[\vspace{-0.3em}\rule{\textwidth}{0.4pt}]
\titlespacing*{\section}{0pt}{0.4em}{0.15em}

\newcommand{\blueitem}[1]{\textcolor{subtitleblue}{\textbf{#1}}}
\newcommand{\sub}[1]{\textcolor{subtle}{\footnotesize #1}}

\setlist[itemize]{label=\textbullet, leftmargin=*, noitemsep, topsep=0pt, parsep=0pt}

\hypersetup{colorlinks=true, urlcolor=linkblue, linkcolor=linkblue}

%----------------------------------------------------------------------------------------
%	DOCUMENT CONTENT
%----------------------------------------------------------------------------------------
\begin{document}

%--- HEADER ---
{\Large \textbf{\textcolor{titleblue}{Iván Sevillano-García}}} \hfill {\footnotesize \textcolor{subtle}{521 citas \textbar\ índice h 4 \textbar\ 512 desde 2021}} \\[0.1em]
{\footnotesize Instituto DaSCI, Universidad de Granada, España} \\[0.1em]
{\footnotesize
\faEnvelope\ \href{mailto:isevillano@ugr.es}{isevillano@ugr.es} \quad
\faGlobe\ \href{https://www.linkedin.com/in/iván-sevillano-garcía-304801116}{LinkedIn} \quad
\faGithub\ \href{https://github.com/isega24}{GitHub} \quad
\faGraduationCap\ \href{https://scholar.google.com/citations?user=g5xyauEAAAAJ}{Scholar} \quad
\faOrcid\ \href{https://orcid.org/0000-0002-5029-9106}{ORCID}
}

\vspace{0.2em}

%--- PROFILE ---
\section*{Perfil}

Investigador en Inteligencia Artificial centrado en IA explicable, detección de datos fuera de distribución y despliegue confiable en ámbitos de alto riesgo. Desarrollo y analizo métodos que mejoran la interpretabilidad, la robustez y la seguridad de los sistemas de aprendizaje profundo, combinando modelado teórico con evaluación empírica sobre conjuntos de datos de referencia y reales. Mis contribuciones abarcan la auditoría de explicaciones, la regularización guiada por explicaciones y la detección OOD explicable con confianza estadística --- avanzando hacia sistemas de IA auditables por diseño y responsables ante las personas a las que afectan. Publicaciones en \textit{Pattern Recognition}, \textit{Machine Intelligence Research} e \textit{IEEE JBHI}; ponente en la Escuela Internacional de Verano de ADIA Lab sobre XAI (2025); mantengo activamente librerías XAI y conjuntos de datos de referencia en código abierto. Cultivo la investigación mediante el diálogo informal y continuo con colegas de otras instituciones: el marco STOOD-X, por ejemplo, nació de una propuesta de grupo de trabajo en la Escuela de Verano de ADIA Lab (2024) y maduró hasta una publicación en \textit{Pattern Recognition}. Actualmente extiendo la IA confiable hacia sistemas agénticos ---ingeniería de \textit{harness}, ingeniería de bucles y memoria persistente con LLM de última generación--- con el objetivo futuro de una investigación interdisciplinar que conecte aprendizaje automático, estadística y conocimiento de dominio en salud, finanzas y ciberseguridad. Mi meta es crecer hacia el liderazgo de equipos de investigación, formando grupos que entreguen IA confiable con impacto medible en el mundo real mediante mentoría, coordinación y dirección interdisciplinar.

%--- RESEARCH INTERESTS ---
\section*{Intereses de investigación}

IA explicable (XAI) \quad\textbar\quad Detección de datos fuera de distribución \quad\textbar\quad Evaluación y auditoría de explicaciones \quad\textbar\quad Regularización guiada por explicaciones \quad\textbar\quad IA agéntica confiable \quad\textbar\quad Ingeniería de \textit{harness} y de bucles

%--- EDUCATION ---
\section*{Formación}

\blueitem{Doctorado en Tecnologías de la Información y la Comunicación \hfill 2021--2025} \\
\sub{Universidad de Granada --- Instituto DaSCI} \\
\textit{Tesis:} Explainable AI: Methods for the Generation, Evaluation, and Correction of Explanations \\
\sub{Directores: F.~Herrera, J.~Luengo \quad Defensa: 13 nov 2025 \quad \textbf{Sobresaliente Cum Laude}} \\
\sub{Contribuciones: REVEL (evaluación de explicaciones), X-SHIELD (regularización de explicaciones), STOOD-X (detección OOD explicable).} \\
\sub{\href{https://hdl.handle.net/10481/112935}{hdl.handle.net/10481/112935} \quad ISBN 9791370150433}

\vspace{0.2em}

\blueitem{Máster en Ciencia de Datos e Ingeniería de Computadores \hfill 2019--2020} \\
\sub{Universidad de Granada}

\vspace{0.2em}

\blueitem{Doble Grado en Ingeniería Informática y Matemáticas \hfill 2013--2019} \\
\sub{Universidad de Granada --- ETSIIT}

%--- RESEARCH EXPERIENCE ---
\section*{Experiencia investigadora}

\blueitem{Investigador predoctoral \hfill 2019--2025} \\
\sub{Universidad de Granada --- Instituto DaSCI}
\begin{itemize}
    \item \textbf{REVEL:} diseño de un marco teórico para auditar explicaciones lineales locales en robustez, concisión y prescriptividad; evaluado sobre \textit{benchmarks} de clasificación de imágenes con aprendizaje profundo.
    \item \textbf{X-SHIELD:} propuesta de una familia de regularizadores guiados por explicaciones que seleccionan características de entrada que ocultar durante el entrenamiento, mejorando simultáneamente rendimiento y explicabilidad del modelo; publicado en \textit{Machine Intelligence Research}.
    \item \textbf{STOOD-X:} desarrollo de un método de detección OOD explicable en dos etapas que combina el test no paramétrico de Wilcoxon-Mann-Whitney con explicaciones visuales basadas en conceptos; publicado en \textit{Pattern Recognition}.
    \item \textbf{COVIDGR:} codiseño de un conjunto de radiografías de tórax COVID-19 equilibrado en todos los niveles de gravedad y de la metodología COVID-SDNet; publicado en \textit{IEEE JBHI} (477 citas).
    \item Investigador en la Cátedra ENIA IAFER (programa financiado por EU Next Generation sobre IA responsable).
\end{itemize}

\blueitem{Investigador visitante \hfill feb--mar 2026} \\
\sub{Universidad Politécnica de Poznań, Polonia --- anfitrión: Prof.~J.~Stefanowski}
\begin{itemize}
    \item Investigación sobre explicaciones basadas en prototipos y contrafactuales para métodos XAI post-hoc; propuesta de publicación conjunta en curso.
\end{itemize}

\blueitem{Colaboración industrial --- ArcelorMittal \hfill 2020--2021} \\
\sub{vía DaSCI, UGR}
\begin{itemize}
    \item Codesarrollo de un \textit{ensemble} de aprendizaje profundo para la segmentación de imágenes metalográficas; publicación conjunta del conjunto de datos MetalDAM; publicado en \textit{Information Fusion} (59 citas).
\end{itemize}

%--- PUBLICATIONS ---
\section*{Publicaciones}

\footnotesize

\blueitem{Revistas}

\textbf{STOOD-X: Explainable OOD Detection via Nonparametric Statistical Testing on Large-Scale Datasets.} \\
\sub{\textit{Pattern Recognition}, 2026. I.~Sevillano-García, J.~Luengo, F.~Herrera. \href{https://doi.org/10.1016/j.patcog.2026.113254}{DOI}}

\textbf{X-SHIELD: Regularization for Explainable Artificial Intelligence.} \\
\sub{\textit{Machine Intelligence Research}, 23, 528--544, 2026. I.~Sevillano-García, J.~Luengo, F.~Herrera. \href{https://doi.org/10.1007/s11633-025-1576-y}{DOI}}

\textbf{A Three-Level Framework for LLM-Enhanced eXplainable AI.} \\
\sub{\textit{Information Systems Frontiers}, 2025. M.~Bello, R.~Bello, M.M.~García, A.~Nowé, I.~Sevillano-García, F.~Herrera. \href{https://doi.org/10.1007/s10796-025-10668-1}{DOI}}

\textbf{REVEL Framework to Measure Local Linear Explanations for Black-Box Models.} \\
\sub{\textit{Intl.~J.~Intelligent Systems}, 2023. I.~Sevillano-García, J.~Luengo, F.~Herrera. \href{https://doi.org/10.1155/2023/8068569}{DOI} \quad (15 citas)}

\textbf{A Tutorial on Segmentation of Metallographic Images: MetalDAM Dataset, DL Ensemble.} \\
\sub{\textit{Information Fusion}, 2022. J.~Luengo, R.~Moreno, I.~Sevillano-García, et~al. \href{https://doi.org/10.1016/j.inffus.2021.09.018}{DOI} \quad (59 citas)}

\textbf{COVIDGR Dataset \& COVID-SDNet Methodology for Predicting COVID-19 (Chest X-ray).} \\
\sub{\textit{IEEE JBHI}, 24(12), 3595--3605, 2020. S.~Tabik, A.~Gómez-Ríos, J.L.~Martín-Rodríguez, I.~Sevillano-García, et~al. \href{https://doi.org/10.1109/JBHI.2020.3037127}{DOI} \quad (477 citas)}

\blueitem{Congresos}

\textbf{Fine-Grained Prostate Cancer Tissue Classification via Expert-Annotated Dataset \& Interpretable DL.} \\
\sub{\textit{IEEE CAI 2026}, 1264--1269. J.~Ayensa-Jiménez, D.~Navarro, I.~Sevillano-García, et~al. \href{https://doi.org/10.1109/CAI68641.2026.11536250}{DOI}}

\textbf{Optimizing LIME Explanations Using REVEL Metrics.} \\
\sub{\textit{HAIS}, 304--313, 2023. I.~Sevillano-Garcia, J.~Luengo, F.~Herrera. \href{https://doi.org/10.1007/978-3-031-40725-3_26}{DOI}}

\textbf{Low-Impact Feature Reduction Regularization Term: How to Improve AI with Explainability.} \\
\sub{\textit{xAI 2023 (Late-breaking Work, poster)}, 135--139. I.~Sevillano-García, J.~Luengo, F.~Herrera. \href{https://dblp.org/rec/conf/xai/Sevillano-Garcia23}{DBLP}}

\textbf{Slicer: Feature Learning for Class Separability with LS-SVM Loss (COVID-19 X-ray).} \\
\sub{\textit{HAIS}, 305--315, 2021. D.~Charte, I.~Sevillano-García, et~al. \href{https://doi.org/10.1007/978-3-030-86271-8_26}{DOI}}

\normalsize

%--- OPEN SOURCE ---
\section*{Software y datos de código abierto}

\begin{itemize}
    \item \textbf{S-ReVel} (\href{https://github.com/ari-dasci/S-ReVel}{github.com/ari-dasci/S-ReVel}) --- librería de auditoría XAI, mantenida activamente. AGPL-3.0.
    \item \textbf{S-STOOD-X} (\href{https://github.com/ari-dasci/S-STOOD-X}{github.com/ari-dasci/S-STOOD-X}) --- librería de detección OOD explicable, mantenida activamente. GPL-3.0.
    \item \textbf{S-T-SHIELD} (\href{https://github.com/ari-dasci/S-T-SHIELD}{github.com/ari-dasci/S-T-SHIELD}) --- familia de regularizadores T-SHIELD, mantenida activamente. GPL-3.0.
    \item \textbf{COVIDGR} --- conjunto de datos de referencia de radiografías de tórax COVID-19 (\textit{IEEE JBHI} 2020; 477 citas).
    \item \textbf{MetalDAM} --- conjunto de imágenes metalográficas (\textit{Information Fusion} 2021; 59 citas).
\end{itemize}

%--- TEACHING ---
\section*{Docencia y tutorización}

\begin{itemize}
    \item Codirección de 2 doctorandos en la UCLV (Cuba), en colaboración con el grupo del Prof.~M.~Bello --- en curso.
    \item Codirección de doctorandos (en curso): \textit{Machine Unlearning}, métricas XAI, XAI para agentes de IA.
    \item Mentoría continuada de un investigador doctoral en un marco general de auditoría de sistemas de IA.
    \item Codirección de trabajos de fin de grado y de máster, UGR.
    \item Docencia universitaria, Dpto.~de Ciencias de la Computación e IA, UGR.
\end{itemize}

%--- COLLABORATION & RESEARCH DIALOGUE ---
\section*{Colaboración y diálogo investigador}

\begin{itemize}
    \item Mantengo una práctica informal de discusión de ideas de investigación y líneas de proyecto con colegas de otras instituciones, sembrando colaboraciones más allá de las estructuras formales de dirección de tesis.
    \item Coorganicé la propuesta de un grupo de trabajo sobre detección OOD en la Escuela Internacional de Verano de ADIA Lab (Granada, 2024) junto con otros asistentes --- la semilla del marco STOOD-X, publicado después en \textit{Pattern Recognition} (2026).
    \item Lidero sesiones de lluvia de ideas del grupo para la ideación de investigación y la revisión entre pares.
\end{itemize}

%--- ACADEMIC SERVICE & FUNDING ---
\section*{Actividad académica y financiación}

\blueitem{Ponencias invitadas:} Escuela Internacional de Verano de ADIA Lab sobre IA Explicable --- ``Lab on xAI Tools'' (Granada, mayo 2025); organicé y coordiné las actividades de los equipos de trabajo con los estudiantes participantes. \\
\blueitem{Revisión por pares:} \textit{Information Sciences} (2025), \textit{Information Fusion} (2023). Verificado vía ORCID. \\
\blueitem{Proyectos:} Cátedra ENIA IAFER (EU NextGen, TSI-100927-2023-1); IAFER CiB --- Cátedra INCIBE de IA Ética para Ciberseguridad (DaSCI/UGR); PID2023-150070NB-I00 (MCIN/AEI).

%--- TECHNICAL SKILLS ---
\section*{Competencias técnicas}

\blueitem{Lenguajes:} Python, R \\
\blueitem{Entornos:} PyTorch, TensorFlow, scikit-learn \\
\blueitem{Herramientas:} Git, Linux, Conda, TensorBoard, Slurm \\
\blueitem{Computación:} clústeres multi-GPU (NVIDIA A100), entrenamiento distribuido, seguimiento de experimentos

%--- LANGUAGES ---
\section*{Idiomas}

Español --- nativo \quad\textbar\quad Inglés --- fluido (competencia profesional plena; ponencias impartidas en inglés)

%--- REFERENCES ---
\section*{Referencias}

\begin{minipage}[t]{0.48\textwidth}
    \textbf{Prof.~Francisco Herrera} \\
    \sub{Catedrático de Universidad \\ Dpto.~de CC.~de la Computación e IA, UGR \\ \href{mailto:herrera@decsai.ugr.es}{herrera@decsai.ugr.es}}
\end{minipage}\hfill
\begin{minipage}[t]{0.48\textwidth}
    \textbf{Prof.~Julián Luengo} \\
    \sub{Catedrático de Universidad \\ Dpto.~de CC.~de la Computación e IA, UGR \\ \href{mailto:julianlm@decsai.ugr.es}{julianlm@decsai.ugr.es}}
\end{minipage}

\end{document}
